% Standalone resolution manuscript generated from solution.md.
% Compile with XeLaTeX; author and review metadata are maintained in Markdown.
\documentclass[10pt,a4paper]{article}
\usepackage[margin=24mm,headheight=15pt]{geometry}
\usepackage{fontspec}
\setmainfont{DejaVuSerif.ttf}[BoldFont=DejaVuSerif-Bold.ttf,ItalicFont=DejaVuSerif-Italic.ttf,BoldItalicFont=DejaVuSerif-BoldItalic.ttf]
\setsansfont{DejaVuSans.ttf}[BoldFont=DejaVuSans-Bold.ttf,ItalicFont=DejaVuSans-Oblique.ttf]
\setmonofont{DejaVuSansMono.ttf}
\usepackage{amsmath,amssymb,unicode-math}
\setmathfont{latinmodern-math.otf}
\usepackage{xcolor,microtype,parskip,needspace,fancyhdr}
\usepackage{bookmark}
\usepackage{xurl}
\hypersetup{colorlinks=true,urlcolor=blue,linkcolor=blue,pdftitle={SP-11:
The delta bound from Hall's theorem},pdfauthor={George
Stepaniants},pdflang={en-GB}}
\urlstyle{same}
\setlength{\emergencystretch}{3em}
\providecommand{\tightlist}{\setlength{\itemsep}{0pt}\setlength{\parskip}{0pt}}
\providecommand{\pandocbounded}[1]{#1}
\setcounter{secnumdepth}{0}
\allowdisplaybreaks[1]
\widowpenalty=10000
\clubpenalty=10000
\pagestyle{fancy}
\fancyhf{}
\fancyhead[L]{\small\sffamily APPLICATION NOTE / SP-11}
\fancyhead[R]{\small\sffamily 11 September 2026}
\fancyfoot[L]{\parbox[b]{0.9\textwidth}{\fontsize{8}{10}\selectfont Hall's
theorem; AI-assisted application note and independent automated-agent
review.}}
\fancyfoot[R]{\footnotesize\thepage}
\begin{document}
{\raggedright\Large\bfseries SP-11: The delta bound from Hall's
theorem\par}
\vspace{0.4em}
{\normalsize\bfseries George Stepaniants\par}
{\small Department of Computing and Mathematical Sciences, California
Institute of Technology, Pasadena, California, USA.\par}
\vspace{0.6em}
\pagestyle{plain}

\textbf{Literature-dependent resolution.} The all-graph theorem used
here is due to H. Tracy Hall. This application note was prepared with
substantial ChatGPT/Codex assistance. A separate
\href{https://github.com/ajt60gaibb/OpenProblemsInNLA/blob/main/references/stepaniants-sp11-sp12-2026-09-11/verification/SP-11-SP-12-independent-review.md}{independent
Codex-agent review} checked this deduction and the essential proof of
Hall's theorem and returned \textbf{PASS}. This is automated-agent
mathematical review, not external human peer review, formal
verification, or a claim that the cited preprint has been refereed.

\subsection{Statement and notation}\label{statement-and-notation}

Let \(G\) be a finite simple undirected graph on \(n\ge1\) vertices, and
let \(\mathcal S(G)\) be the real symmetric \(n\times n\) matrices whose
off-diagonal nonzero entries occur exactly at the edges of \(G\).
Diagonal entries are unrestricted. Writing \(\delta(G)\) for the minimum
vertex degree, the target is

\[
\operatorname{mr}(G):=\min_{A\in\mathcal S(G)}\operatorname{rank}A
\le n-\delta(G).
\]

Define \(\nu(G)\) as the maximum nullity among positive-semidefinite
\(A\in\mathcal S(G)\) satisfying the strong Arnold property (SAP): the
only real symmetric \(X\) with

\[
AX=0,\qquad A\circ X=0,\qquad I_n\circ X=0
\]

is \(X=0\), where \(\circ\) denotes entrywise product.

\subsection{External theorem}\label{external-theorem}

Hall's \textbf{Corollary 3.22}, based on \textbf{Theorem 3.20}, proves

\[\nu(G)\ge\delta(G)\]

for every finite simple graph {[}H{]}. The source is arXiv:2601.01211v1,
submitted 3 January 2026; the current record still listed only v1 when
checked on 11 September 2026. The theorem is used as an external input
here. The independent review records its full essential proof audit and
the nonessential wording issues in the source. In particular, the matrix
in the definitions and Gram construction is positive
\textbf{semidefinite}, despite the word ``definite'' in the abstract.

\subsection{Deduction: affirmative answer to
SP-11}\label{deduction-affirmative-answer-to-sp-11}

Choose a positive-semidefinite SAP matrix \(A\in\mathcal S(G)\) with
nullity at least \(\delta(G)\), as provided by {[}H{]}. Forgetting the
extra PSD and SAP conditions leaves an admissible matrix for SP-11.
Rank-nullity gives

\[
\operatorname{mr}(G)\le\operatorname{rank}A
=n-\operatorname{nullity}A\le n-\delta(G).
\]

This is exactly the requested inequality. No connectedness assumption,
prescribed edge weights, or diagonal restriction is introduced. When
\(\delta(G)=0\), the bound also holds trivially. \(\square\)

\subsection{Attribution and scope}\label{attribution-and-scope}

The substantive all-graph existence theorem belongs to Hall. His
\textbf{Corollary 3.24} explicitly states the ordinary Delta Conjecture
as a consequence. George Stepaniants is the author of this explanatory
application note, not the discoverer of that theorem. No novelty or
priority claim is made.

The
\href{https://github.com/ajt60gaibb/OpenProblemsInNLA/blob/main/references/stepaniants-sp11-sp12-2026-09-11/README.md}{submission
record} preserves the original reconstructed note, the independent
review, and the dated public-network audit. Unsupported earlier claims
about unavailable artifacts and graph-atlas computations remain
withdrawn; they are not evidence for this resolution.

\subsection{References}\label{references}

{[}H{]} H. Tracy Hall, \emph{The Delta Theorem: A dimension bound for
faithful orthogonal graph representations}, arXiv:2601.01211v1, 3
January 2026, Theorem 3.20 and Corollaries 3.22-3.24.
\href{https://arxiv.org/html/2601.01211v1}{Versioned primary text} ·
\href{https://arxiv.org/abs/2601.01211}{Current arXiv record}.

{[}R11{]} \emph{Open Problems in Numerical Linear Algebra},
\href{https://github.com/ajt60gaibb/OpenProblemsInNLA/blob/main/eigenvalues-and-inverse-problems/SP-11/README.md}{SP-11:
original canonical target}.
\end{document}
